\begin{lemma}[Uniqueness of the closure vertex set]
If $T_1$ and $T_2$ are two closures of the same tree-sequence $T$ under
the same coloring $\varphi$, then $V(T_1)=V(T_2)$.
\end{lemma}
\begin{proof}
We first prove the following containment principle.  Let $C$ be a closure
of $T$, and let $S$ be any tree-sequence reached from $T$ by Tashkinov
augmentations under the same coloring.  Then $V(S)\subseteq V(C)$.

Because $C$ is reached from $T$, repeated use of the initial-segment clause
in \noderef{TashkinovAugmentation} shows that $V(T)\subseteq V(C)$.  Induct
on the augmentation sequence from $T$ to $S$ supplied by
\noderef{TashkinovAugmentationSequence}.  The empty sequence is the
containment just noted.  For the induction step, suppose
$S^+=(S,f,y)$ is obtained by one augmentation from $S$ and
$V(S)\subseteq V(C)$.  The new edge identity $f$ has an end $x\in V(S)$,
its other end is the new vertex $y\notin V(S)$, and
$\varphi(f)=c$ for a color $c$ missing at some $z\in V(S)$, by
\noderef{TashkinovAugmentation}.  The induction hypothesis puts both $x$
and $z$ in $V(C)$.  If $y\notin V(C)$, then $f$ is a boundary edge identity
of $V(C)$ by \noderef{BoundaryEdges}, colored $c$ while $c$ is missing at
the vertex $z\in V(C)$.  This contradicts the closedness clause in
\noderef{TashkinovClosure}.  Consequently $y\in V(C)$, and hence
$V(S^+)\subseteq V(C)$.  This completes the induction.

Apply the containment principle first with $C=T_2$ and $S=T_1$, using the
augmentation sequence contained in the closure assertion for $T_1$.  It
gives $V(T_1)\subseteq V(T_2)$.  Interchanging $T_1,T_2$ gives the reverse
containment.  Extensionality of finite vertex sets yields
$V(T_1)=V(T_2)$.
\end{proof}
